\documentclass[12pt,a4paper]{article}
\usepackage{../../../myconfig}

\begin{document}

% =========================================================================
% BÀI TẬP VỀ NHÀ: XÁC SUẤT CƠ BẢN VÀ XÁC SUẤT CỔ ĐIỂN (20 CÂU)
% Câu 1-10: các câu/phần chưa dùng trong BTDD, từ tài liệu gốc.
% Câu 11-20: câu bổ sung.
% Bảng đáp án và ghi chú kiểm tra đặt ở cuối file, dưới dạng comment.
% =========================================================================
\titlebox{Xác suất}{Xác suất cơ bản \\ Xác suất cổ điển}{Xác suất: Xác suất cơ bản và xác suất cổ điển}

\vspace{-0.25cm}

\begin{center}
    \textit{Các câu hỏi xác suất: ghi số thập phân, làm tròn đến hàng phần trăm.}
\end{center}

% --- Câu 1 (Câu 4a trong tài liệu gốc) ---
\questionShortAnswer{0}{1}{%
    Một công ty truyền thông đấu thầu hai dự án.
    Xác suất thắng dự án 1 và dự án 2 lần lượt là $50\%$ và $60\%$;
    xác suất thắng cả hai dự án là $40\%$.
    Tính xác suất công ty thắng đúng một dự án.
}


% --- Câu 2 (Câu 5 trong tài liệu gốc, làm rõ cách hỏi) ---
\questionShortAnswer{0}{2}{%
    Trong một vùng, tỉ lệ người hút thuốc và mắc bệnh ung thư phổi là $15\%$;
    hút thuốc và không mắc bệnh là $25\%$;
    không hút thuốc và không mắc bệnh là $50\%$;
    không hút thuốc và mắc bệnh là $10\%$.
    Tỉ lệ mắc bệnh trong nhóm hút thuốc gấp bao nhiêu lần tỉ lệ mắc bệnh trong nhóm không hút thuốc?
}


% --- Câu 3 (Câu 13 trong tài liệu gốc, bổ sung không hoàn lại) ---
\questionShortAnswer{0}{3}{%
    Một thùng có 50 thẻ giảm giá cùng kích thước và chất liệu,
    gồm 30 thẻ xanh và 20 thẻ trắng.
    Lấy ngẫu nhiên lần lượt hai thẻ, không hoàn lại.
    Tính xác suất lần thứ nhất lấy được thẻ xanh và lần thứ hai lấy được thẻ trắng.
}


% --- Câu 4 (Câu 2 trong tài liệu gốc) ---
\questionShortAnswer{0}{4}{%
    Quan sát hai mã cổ phiếu $X$ và $Y$ trong một phiên giao dịch.
    Nếu $Y$ không giảm giá thì xác suất $X$ giảm giá là $\dfrac35$.
    Nếu $X$ không giảm giá thì xác suất $Y$ giảm giá là $\dfrac23$.
    Xác suất cả hai cùng giảm giá là $\dfrac1{10}$.
    Tính xác suất có ít nhất một mã cổ phiếu giảm giá.
}


% --- Câu 5 (Câu 9 trong tài liệu gốc) ---
\questionShortAnswer{0}{5}{%
    Trong cùng một ngày, nếu xã $B$ không mưa thì xác suất xã $A$ không mưa là $65\%$;
    nếu xã $A$ không mưa thì xác suất xã $B$ không mưa là $60\%$.
    Xác suất cả hai xã cùng có mưa là $10\%$.
    Tính xác suất có ít nhất một trong hai xã có mưa.
}


\newpage

% --- Câu 6 (Câu 10 trong tài liệu gốc) ---
\questionShortAnswer{0}{6}{%
    Quan sát hai mã cổ phiếu $A$ và $B$ trong một phiên giao dịch.
    Nếu $B$ không giảm giá thì xác suất $A$ giảm giá là $\dfrac25$;
    nếu $A$ không giảm giá thì xác suất $B$ giảm giá là $\dfrac47$.
    Xác suất cả hai cùng giảm giá là $0{,}1$.
    Tính xác suất có ít nhất một trong hai mã cổ phiếu giảm giá.
}


% --- Câu 7 (Câu 16 trong tài liệu gốc) ---
\questionShortAnswer{0}{7}{%
    Có 9 quả cầu cùng kích thước và khối lượng, được đánh số từ 1 đến 9.
    Lấy ngẫu nhiên lần lượt 4 quả cầu, không hoàn lại.
    Tính xác suất tổng bốn số ghi trên các quả cầu lấy được bằng 15.
}


% --- Câu 8 (Câu 14 trong tài liệu gốc) ---
\questionShortAnswer{0}{8}{%
    Cho $X=\{1;2;\ldots;12\}$.
    Chọn đều ngẫu nhiên 4 số phân biệt từ $X$,
    rồi chọn đều một trong 4 số đặt vào vòng tròn lớn ở giữa;
    ba số còn lại đặt vào ba vòng tròn nhỏ không phân biệt vị trí.
    Gọi $P$ là xác suất tổng hai số ở hai vòng tròn nhỏ bất kì
    đều nhỏ hơn số ở giữa, đồng thời tổng cả ba số ở ngoài lớn hơn số ở giữa.
    Tính $1980P$.
}


% --- Câu 9 (Câu 15 trong tài liệu gốc) ---
\questionShortAnswer{0}{9}{%
    Một hộp có 25 lá thăm khác nhau, được đánh số từ 1 đến 25.
    Chọn ngẫu nhiên đồng thời 4 lá thăm.
    Tính xác suất có ít nhất 2 lá thăm mang số nguyên tố
    và tổng bốn số được chọn là số chẵn.
}


% --- Câu 10 (Câu 19 trong tài liệu gốc, làm rõ mô hình đặt mìn) ---
\questionShortAnswer{0}{10}{%
    Một bảng $9\times9$ có 27 ô chứa mìn;
    tập 27 ô này được chọn đều ngẫu nhiên trong tất cả các tập 27 ô của bảng.
    Đánh số hàng và cột từ 1 đến 9.
    Mỗi ô an toàn hiển thị số mìn trong các ô kề chung cạnh hoặc chung đỉnh với nó.
    Hai ô $(4;5)$ và $(6;5)$ đã mở, đều an toàn,
    lần lượt hiển thị các số 4 và 1.
    Tính xác suất ô trung tâm $(5;5)$ chứa mìn khi biết hai thông tin trên.
}


% --- Câu 11 (Bổ sung) ---
\questionShortAnswer{0}{11}{%
    Một hộp có 6 bi đỏ, 4 bi xanh và 3 bi vàng, các viên bi cùng kích thước và khối lượng.
    Chọn ngẫu nhiên đồng thời 5 viên bi.
    Tính xác suất nhóm được chọn có đúng 2 bi đỏ và có đủ cả ba màu.
}


% --- Câu 12 (Bổ sung) ---
\questionShortAnswer{0}{12}{%
    Gọi $S$ là tập tất cả các số tự nhiên có bốn chữ số đôi một khác nhau,
    được lập từ các chữ số $0;1;2;3;4;5;6;7$.
    Chọn đều ngẫu nhiên một số trong $S$.
    Tính xác suất số được chọn chia hết cho 5.
}


% --- Câu 13 (Bổ sung) ---
\questionShortAnswer{0}{13}{%
    Chọn ngẫu nhiên đồng thời 4 thẻ từ 12 thẻ được đánh số từ 1 đến 12.
    Tính xác suất có ít nhất 2 thẻ mang số chia hết cho 3
    và tổng bốn số được chọn là số chẵn.
}


% --- Câu 14 (Bổ sung) ---
\questionShortAnswer{0}{14}{%
    Một hộp có 3 bi đỏ, 5 bi trắng và 2 bi xanh, các viên bi cùng kích thước và khối lượng.
    Chọn ngẫu nhiên đồng thời 3 viên bi.
    Biết nhóm được chọn có ít nhất một bi đỏ,
    tính xác suất nhóm đó có đúng 2 bi đỏ.
}


% --- Câu 15 (Bổ sung) ---
\questionShortAnswer{0}{15}{%
    Ba học sinh làm bài độc lập với nhau.
    Xác suất làm đúng của từng học sinh lần lượt là $0{,}8$; $0{,}7$ và $0{,}6$.
    Biết có ít nhất một học sinh làm đúng,
    tính xác suất có đúng hai học sinh làm đúng.
}


% --- Câu 16 (Bổ sung) ---
\questionShortAnswer{0}{16}{%
    Có 10 chìa khóa ứng với 10 ngôi nhà, mỗi chìa chỉ mở được đúng một nhà.
    Có 4 nhà không khóa cửa.
    Một người chọn đều ngẫu nhiên 3 chìa mang theo,
    rồi chọn đều ngẫu nhiên một trong 10 nhà, độc lập với việc chọn chìa.
    Tính xác suất người đó vào được nhà đã chọn.
}


% --- Câu 17 (Bổ sung) ---
\questionShortAnswer{0}{17}{%
    Một nhóm có 8 học sinh nam và 6 học sinh nữ.
    Chọn ngẫu nhiên một đội gồm 5 học sinh.
    Tính xác suất đội có ít nhất 2 nam và ít nhất 2 nữ.
}


% --- Câu 18 (Bổ sung) ---
\questionShortAnswer{0}{18}{%
    Xếp ngẫu nhiên 6 học sinh $A$, $B$, $C$, $D$, $E$, $F$ thành một hàng ngang,
    các cách xếp có cùng khả năng.
    Tính xác suất $A$, $B$ đứng cạnh nhau nhưng $C$, $D$ không đứng cạnh nhau.
}


% --- Câu 19 (Bổ sung) ---
\questionShortAnswer{0}{19}{%
    Một hộp có 20 thẻ được đánh số từ 1 đến 20.
    Chọn ngẫu nhiên đồng thời 3 thẻ.
    Tính xác suất tích ba số ghi trên các thẻ được chọn chia hết cho 6.
}


% --- Câu 20 (Bổ sung) ---
\questionShortAnswer{0}{20}{%
    Xếp ngẫu nhiên 3 học sinh nữ và 4 học sinh nam thành một hàng ngang,
    các cách xếp có cùng khả năng.
    Tính xác suất không có hai học sinh nữ nào đứng cạnh nhau.
}


% =========================================================================
% BẢNG ĐÁP ÁN - CHỈ HIỂN THỊ TRONG MÃ NGUỒN
% =========================================================================
% Câu     |  1    |  2    |  3    |  4    |  5
% Đáp án  | 0,30  | 2,25  | 0,24  | 0,80  | 0,59
%
% Câu     |  6    |  7    |  8    |  9    | 10
% Đáp án  | 0,70  | 0,05  | 16    | 0,23  | 0,29
%
% Câu     | 11    | 12    | 13    | 14    | 15
% Đáp án  | 0,35  | 0,27  | 0,21  | 0,25  | 0,46
%
% Câu     | 16    | 17    | 18    | 19    | 20
% Đáp án  | 0,58  | 0,70  | 0,20  | 0,61  | 0,29
%
% =========================================================================
% KIỂM TRA ĐÁP ÁN (GIỮ GIÁ TRỊ CHÍNH XÁC TRƯỚC KHI LÀM TRÒN)
% =========================================================================
%  1. 0,5 + 0,6 - 2*0,4 = 3/10.
%  2. (15/40)/(10/60) = 9/4. So sánh tỉ lệ bệnh TRONG từng nhóm.
%  3. (30/50)*(20/49) = 12/49.
%  4. Đặt t = P(cả hai không giảm): 1/10 + (3/2)t + 2t + t = 1.
%     P(ít nhất một giảm) = 4/5.
%  5. Đặt t = P(cả hai không mưa): 1/10 + (7/13)t + (2/3)t + t = 1.
%     P(ít nhất một mưa) = 509/860.
%  6. Đặt t = P(cả hai không giảm): 1/10 + (2/3)t + (4/3)t + t = 1.
%     P(ít nhất một giảm) = 7/10.
%  7. Có 6 nhóm: (1,2,3,9), (1,2,4,8), (1,2,5,7),
%     (1,3,4,7), (1,3,5,6), (2,3,4,6). P = 6/C_9^4 = 1/21.
%  8. Gọi m là số giữa; a<b<c là ba số ngoài: b+c<m<a+b+c.
%     Số bộ theo m=8,9,10,11,12 lần lượt là 1,1,3,4,7; tổng 16.
%     Không gian mẫu có 4*C_{12}^4 = 1980 kết quả; 1980P = 16.
%  9. Có 2918 nhóm thuận lợi, C_{25}^4 = 12650; P = 1459/6325.
%     Không có số 2: C_8^2*C_{11}^2 + C_8^2*C_5^2 + C_8^3*C_5^1 + C_8^4.
%     Có số 2: (C_{13}^2 - C_5^2)*C_{11}^1.
% 10. Hai vùng quanh ô 4 và ô 1 có 3 ô chung; mỗi vùng có 4 ô riêng chưa mở.
%     Còn 68 ô ngoài hai vùng. Số mìn trong 3 ô chung chỉ có thể là 0 hoặc 1.
%     W0 = 4*C_{68}^{22}; W1 = 12*C_{68}^{23}.
%     P(tâm có mìn | thông tin) = (W1/3)/(W0+W1) = 2/7.
%     Theo mô hình chọn đều 27 vị trí mìn, đáp án là 0,29, khác key gốc 0,23.
% 11. C_6^2*(C_4^1*C_3^2 + C_4^2*C_3^1)/C_{13}^5 = 50/143.
% 12. Tổng 7*7*6*5 = 1470 số; thuận lợi 7*6*5 + 6*6*5 = 390 số.
%     P = 13/49.
% 13. Liệt kê các nhóm tăng dần: 105 nhóm thuận lợi / C_{12}^4 = 7/33.
%     Chia theo số bội của 3 được chọn: đúng 2 có 88; đúng 3 có 16; đúng 4 có 1.
% 14. C_3^2*C_7^1/(C_{10}^3 - C_7^3) = 21/85.
% 15. Tử = 0,8*0,7*0,4 + 0,8*0,3*0,6 + 0,2*0,7*0,6 = 0,452.
%     Mẫu = 1 - 0,2*0,3*0,4 = 0,976. Tỉ số = 113/244.
% 16. 4/10 + (6/10)*(3/10) = 29/50.
% 17. (C_8^2*C_6^3 + C_8^3*C_6^2)/C_{14}^5 = 100/143.
% 18. (2*P_5 - 4*P_4)/P_6 = 1/5.
% 19. 1 - (C_{10}^3 + C_{14}^3 - C_7^3)/C_{20}^3 = 691/1140.
%     Không chia hết cho 2: chọn toàn số lẻ (10 số).
%     Không chia hết cho 3: chọn toàn số không là bội của 3 (14 số).
%     Giao hai trường hợp: 7 số lẻ không là bội của 3.
% 20. P_4*C_5^3*P_3/P_7 = 2/7.

\end{document}
